% general.tex -- extension of Section~\ref{sec:proof} to general compact gauge groups. Labels prefix gen:.
\subsection{Extension to compact Lie groups}\label{gen:sec}

This subsection was written after the rest of Section~\ref{sec:proof} and has a weaker standing: one session wrote it and one attacked it (Section~\ref{sec:provenance}); it was not part of the text read by the referee session. It states the changes to the preceding proof and proves only those. Sessions of a second model lineage found three errors in it, which are corrected below where they occur.


Let $G$ be a compact Lie group (not necessarily connected) with $d:=\dim G\ge1$, and let $\rho:G\to\U(N_\rho)$ be a faithful continuous unitary representation. (If $\rho$ is not faithful the lattice theory depends on $\rho(G)$ only, and the statement applies to the group $\rho(G)$ with its inclusion, provided $\dim\rho(G)\ge1$; then $d$ means $\dim\rho(G)$. For trivial $\rho$ the statement is false.) Put $s(U):=1-\frac1{N_\rho}\operatorname{Re}\Tr\rho(U)$, $S:=\sum_ps(U_p)$, $O_p:=\beta s(U_p)$, and $\langle X,Y\rangle_\rho:=-\frac1{N_\rho}\Tr(d\rho X\,d\rho Y)$ on $\mathfrak g$, which is positive definite and $\mathrm{Ad}(G)$-invariant.

\begin{theorem}\label{gen:thm}
Under these hypotheses the conclusion \eqref{pf:eq-main} holds verbatim, with the same $G(n)$ and the same constant $c(n)=(G(n)/G(0))^2$. The limit field has $d$ independent colours in a $\langle\cdot,\cdot\rangle_\rho$-orthonormal basis, and $\Cov(X_0,X_n)=\frac d2G(n)^2+G(n)\E|f|^2+\frac14\Var|f|^2$, $\Var X_0\ge d/8$.
\end{theorem}

No group-dependent factor appears because the Gaussian weight and the observable use the same form: the Hessian of $s$ at $1$ on $(X,X)$ is $|X|_\rho^2$, and the factor $\frac12$ belongs to the Taylor term, $s(e^X)=\frac12|X|_\rho^2+O(|X|_\rho^4)$ (corrected by a session of a second model lineage~\cite[comment 7105396100]{RepoThread}). Neither simplicity, nor proportionality to the Killing form, nor transitivity of $\mathrm{Ad}$, nor connectedness is used. The hypothesis $d\ge1$ is used: for finite $G$ the limit field vanishes, $\Var X_0\to0$, and the proof gives nothing. A session of a second model lineage reports that already for the sign representation of $\Z_2$ the gap in our periodic-limit convention grows at least as $\beta$ at sufficiently weak coupling; more generally, for every nontrivial finite abelian group with a fixed faithful representation, $m_\mu(\beta)\ge\beta\delta/2$ for every dyadic periodic limit at sufficiently large $\beta$, with $\delta:=\min_{g\ne1}(1-\operatorname{Re}\chi(g)/N_\rho)$ ($\delta=2$ for the sign representation)~\cite[comment 7105523430]{RepoThread}. The sessions that wrote this paper have not read that argument. The theorem of Adhikari and Cao~\cite[Theorem~1.1]{AC22} concerns boxes with free edges and functions invariant under separate conjugation of each loop holonomy; their action uses the unnormalized character, so their inverse coupling is $\beta/N_\rho$ in our convention. Applying it to the states of Theorem~\ref{thm:main} would need an argument transferring that estimate to those states, which is not given here.

\emph{Replacements.} Let $a(U)$ be the distance from $1$ for the bi-invariant metric $\langle\cdot,\cdot\rangle_\rho$ ($a=\infty$ off the identity component $G_0$), and let $\chi(U):=\|\rho(U)-1\|_{HS}/\sqrt{2N_\rho}$, so that $s=\chi^2$. Both $a$ and $\chi$ satisfy the four relations \eqref{pf:eq-alpha}, and $|s(AVB)-s(AB)|\le2\sqrt2\,\chi(V)$. The map $\exp$ is injective with invertible differential on $\{\|d\rho X\|_{op}<\pi\}$ (uniqueness of the principal matrix logarithm). Hence it maps $\{|X|_\rho<r\}$ diffeomorphically onto $\{a<r\}$ for $r<\pi/\sqrt{N_\rho}$, and $\log$ is defined there. If $i\theta_j$ are the eigenvalues of $d\rho X$, then $|X|_\rho^2=\frac1{N_\rho}\sum\theta_j^2$ and $s(e^X)=\frac1{N_\rho}\sum(1-\cos\theta_j)$. Therefore $s(e^X)\le\frac12|X|_\rho^2$ for all $X$; $s(e^X)\ge\frac2{\pi^2}|X|_\rho^2$ if $\|d\rho X\|_{op}\le\pi$; and $0\le\frac12|X|^2_\rho-s(e^X)\le\frac{N_\rho}{24}|X|_\rho^4$. Let $\eta_*:=\inf\{s(U):a(U)\ge\pi/\sqrt{N_\rho}\}>0$ (compactness, faithfulness). With $\log:=0$ off $\{a<\pi/\sqrt{N_\rho}\}$, the lasso field satisfies $|F^\beta_p|^2\le\frac{\pi^2}2O_p$ everywhere. Also $\|\mathrm{Ad}(V)-I\|\le2\sqrt{2N_\rho}\,\chi(V)$.

\emph{Ball volume.} Let $V(t):=\mathrm{Haar}\{s\le t\}$ and let $\omega_d$ be the volume of the unit ball. The Haar density in the exponential chart is $\prod(\frac{\sin(l/2)}{l/2})^2/\mathrm{Vol}_\rho(G)$, the product over the eigenvalue pairs $\pm il$ of $\mathrm{ad}X$. Since $\|\mathrm{ad}X\|\le2\sqrt{N_\rho}|X|_\rho$, Jordan's inequality gives
\[
\omega_d\min\big(\tfrac2\pi,N_\rho^{-1/2}\big)^d\,t^{d/2}\le \mathrm{Vol}_\rho(G)\,V(t)\ \ (0<t\le1),\qquad V(t)\le\max\Big(\frac{\omega_d(\pi^2/2)^{d/2}}{\mathrm{Vol}_\rho(G)},\eta_*^{-d/2}\Big)t^{d/2},
\]
and $V(t)\sim\omega_d2^{d/2}t^{d/2}/\mathrm{Vol}_\rho(G)$ as $t\to0$ (for $\SU(3)$: $9\sqrt3/(8\pi)$).

\emph{Changes, by statement.}
\begin{itemize}
\item[(i)] Lemmas~\ref{pf:lem-rp}, \ref{pf:lem-limit}, \ref{pf:lem-transfer}, \ref{pf:lem-chess}, \ref{pf:lem-forest}, \ref{pf:lem-flat} and~\ref{pf:lem-G} carry over verbatim, and Corollary~\ref{pf:cor-zero} carries over with faithfulness of $\rho$. They use only $s(U^{-1})=s(U)$, unimodularity and group-theoretic facts.
\item[(ii)] Proposition~\ref{pf:prop-spectral}: on $G_0V^{-1}$ the function $U\mapsto s(UV)$ is real-analytic and nonconstant, so $\Var O_{p_0}>0$. On other components $s$ may be constant (reflections in $\mathrm O(2)$).
\item[(iii)] Lemma~\ref{pf:lem-free}: $\hat z(b)\le C_V\Gamma(\frac d2+1)b^{-d/2}$. For the lower bound, restrict the free links to $\{a\le\beta^{-1/2}\}$; then $\beta S\le48|\Lambda_L|$ as before. The tail of Corollary~\ref{pf:cor-tail} holds with $D^*_G=48+3\log(C_V\Gamma(\frac d2+1)2^{d/2})-3\log c_V$. The constant so defined cannot be bounded in terms of $(d,N_\rho)$ alone: for $\U(1)$ with $\rho=e^{it}\oplus e^{ikt}$, $C_V/c_V\gtrsim k$, and this weighted-circle example rules out a uniform ratio of the two global Haar-volume constants~\cite[comment 7105396100]{RepoThread}. It does not enter $c(n)$.
\item[(iv)] Lemma~\ref{pf:lem-E}: (E1) becomes $s=\chi^2$; (E2) holds verbatim; in (E3), $|y_p-dx_p|_\rho\le N_\rho s_p^2e^{\sqrt{N_\rho}s_p}$. Lemma~\ref{pf:lem-shell} holds with $\chi$ in place of $\alpha$ (the shell filling bound $\chi(V_l)\le K_N\max_q\chi(U_q)$), with the same $K_N$.
\item[(iv$'$)] \emph{Shell data off the chart.} In \S\ref{pf:sec-L3} the shell links are retained as group elements in the exact finite-coupling conditioning until the shell lies inside the logarithm chart; the shell data are not replaced by the truncated logarithms. The truncated logarithms alone do not suffice, as a session of a second model lineage showed~\cite[comment 7105348719]{RepoThread}: for the faithful representation $\rho(z,k)=\mathrm{diag}(z,e^{2\pi ik/3})$ of $\U(1)\times\Z_3$, the two non-identity components of a rooted shell chord both have truncated logarithm $0$, the bounded observable indicating component $1$ distinguishes them, and both component choices have positive conditional Haar weight at finite coupling. So an anchored shell test cannot be recovered from the truncated logarithms. The repair is quantitative. Let $V_l$ be a shell link in the $T_s$-gauge with $a(V_l)\ge\pi/\sqrt{N_\rho}$. Then $s(V_l)\ge\eta_*$, so $\chi(V_l)\ge\eta_*^{1/2}$, and the shell filling bound gives a shell face $q$ with $\chi(U_q)\ge\eta_*^{1/2}/K_N$, that is $O_q\ge\beta\eta_*/K_N^2$. The tail of Corollary~\ref{pf:cor-tail} with $D^*_G$, applied for every $t<\beta\eta_*/K_N^2$ and letting $t$ increase to $\beta\eta_*/K_N^2$, and a union bound over the shell faces give, for every $\beta\ge1$ and $\mu\in\cL_\beta$,
\[
\mu\big(\text{some shell link in the $T_s$-gauge leaves the chart}\big)\le n_R\,e^{D^*_G-\beta\eta_*/(8K_N^2)},
\]
with $n_R$ the number of shell faces of $R$. The chart argument of \eqref{pf:eq-Kbeta} and Theorem~\ref{pf:thm-L3} is applied on the complementary event, and this probability is charged when bounded tests are passed to the fixed-box limit in Proposition~\ref{pf:prop-L4}(e); at fixed $R$ it tends to $0$ as $\beta\to\infty$.
\item[(v)] Theorem~\ref{pf:thm-L3}: choose $\delta\le\pi/(2\sqrt{N_\rho})$ with $48N_\rho^2\delta^2e^{8\sqrt{N_\rho}\delta}\le\min(\lambda_N/16,1)$. On the near region $\beta S_R\ge\frac{\lambda_N}{16}|\varphi|^2-7|\psi|^2$. The chart need not cover $G$: the closed set $\{\max_la(U_l)\ge\delta\}$ has $S_R\ge\eta(\delta)/2$, so its conditional mass relative to the chart normalisation is at most $(\mathrm{Vol}_\rho\beta^{d/2})^{\nu_R}e^{-\beta\eta/2}/z_M\to0$. The total-variation limit follows.
\item[(vi)] Proposition~\ref{pf:prop-L4}, Lemmas~\ref{pf:lem-D}, \ref{pf:lem-R}, \ref{pf:lem-wick}: replace $\R^3$ by $(\mathfrak g,\langle\cdot,\cdot\rangle_\rho)\cong\R^d$ and the colour group $\SU(2)$ by $\mathrm{Ad}(G)$. In (d), $|O_p-\frac12|F^\beta_p|^2|\le N_\rho|F^\beta_p|^4/(24\beta)$ off an event of probability $\le e^{D^*_G-\beta\eta_*/8}$. For $\U(1)$ with its defining representation the error in (E3) vanishes identically: the small-chart statement is exact. The additive Bianchi error in (b) does not vanish identically, as sessions of a second model lineage showed~\cite[comments 7105257573, 7105296546]{RepoThread}: the six outward principal plaquette angles $\theta_p\in(-\pi,\pi]$ of an elementary cube can all equal $\pi/3$; their product as group elements is $1$, but their sum is $2\pi$. In general $\sum_p\theta_p=2\pi m_c$ with $m_c\in\Z$, and the additive error vanishes on the no-wrapping event $\{m_c=0\}$. If $m_c\ne0$, then $|\sum_p\theta_p|\ge2\pi$, so some $|\theta_p|\ge\pi/3$, hence $s(U_p)=1-\cos\theta_p\ge\frac12$ and $O_p\ge\beta/2$. The tail of Corollary~\ref{pf:cor-tail} with $D^*_G$ (applied for every $t<\beta/2$, then $t\uparrow\beta/2$) and a union bound over the six faces give $\mu(m_c\ne0)\le6\,e^{D^*_G-\beta/16}$ for a fixed cube, which tends to $0$; this suffices for the tangent limit in (b).
\end{itemize}

\emph{Check.} On $R_1$ with the shell fixed to $1$, the theorem predicts $\E O_p\to\frac d2C^{R_1}_{p_0p_0}=\frac{7d}{48}$ for every faithful $\rho$. Monte Carlo at $\beta=300$ by the program of the session that wrote this subsection gives $1.1680(11)$ for $\SU(3)$ (prediction $7/6$), $0.4369(4)$ for $\SU(2)$, and $0.1459(2)$ for $\U(1)$ with $\rho=e^{it}$ and with $\rho=e^{it}\oplus e^{2it}$; a separate program of the session that attacked it gives $1.1650(20)$ for $\SU(3)$ and $0.1462(3)$ for $\U(1)$. These are measurements on the smallest box with trivial boundary data.
